\chapter{Pancreas Transplant}

\section*{Introduction \& Clinical Indications}
A pancreas transplant is a complex surgical procedure that involves implanting a healthy pancreas, often with a segment of the donor duodenum, from a deceased or, less commonly, a living donor into a recipient. This intervention is primarily performed to restore endogenous insulin production in individuals, most often those with severe type 1 diabetes mellitus, who suffer from life-threatening complications of their disease. The new pancreas is typically placed heterotopically in the lower abdomen, usually the right iliac fossa, where its endocrine function can physiologically regulate blood glucose levels. This effectively eliminates the need for exogenous insulin injections and can halt or potentially reverse the progression of debilitating diabetic complications, significantly improving both quality of life and long-term survival.

\begin{itemize}
  \item Pancreatic Transplantation
  \item Simultaneous Pancreas-Kidney (SPK) Transplant
  \item Pancreas After Kidney (PAK) Transplant
  \item Pancreas Transplant Alone (PTA)
  \item Whole Pancreas Transplant
\end{itemize}

The fundamental surgical principle of pancreas transplantation is to provide a functional endocrine organ capable of producing insulin in a glucose-responsive manner, thereby achieving sustained euglycemia without exogenous insulin administration. The transplanted pancreas contains the islets of Langerhans, which are clusters of cells responsible for hormone production, including insulin and glucagon. By implanting a healthy, insulin-producing pancreas, the procedure directly addresses the underlying pathophysiology of type 1 diabetes mellitus, which is the autoimmune destruction of the native pancreatic beta cells.

The technical goals involve establishing robust arterial inflow and venous outflow for the donor pancreas, typically by connecting its vascular pedicle to the recipient's iliac vessels. Equally crucial is the management of the exocrine secretions, primarily digestive enzymes. Modern practice favors enteric drainage, where the donor duodenal segment is anastomosed to the recipient's small intestine. This allows the exocrine enzymes to drain physiologically into the gastrointestinal tract, minimizing metabolic complications and providing a more natural environment for the graft. The successful establishment of vascular flow and exocrine drainage allows the transplanted islets to function, releasing insulin directly into the systemic circulation in response to glucose fluctuations, thus restoring metabolic control and preventing or reversing the severe complications of diabetes.

The journey of pancreas transplantation began in the mid-20th century, marked by significant surgical and immunological advancements. The first human pancreas transplant was performed in December 1966 by Richard C. Lillehei and William Kelly at the University of Minnesota. This pioneering procedure involved a segmental pancreatic graft with duodenal drainage into the jejunum, though long-term success was limited due primarily to formidable challenges with organ rejection and surgical complications, particularly graft thrombosis and anastomotic leaks.

Progress remained slow through the 1970s, with poor graft and patient survival rates. A major turning point arrived in the early 1980s with the introduction of cyclosporine, a potent calcineurin inhibitor, which revolutionized immunosuppressive therapy and dramatically improved transplant outcomes. This era also saw the increasing adoption of simultaneous pancreas-kidney (SPK) transplantation, particularly for patients with type 1 diabetes and end-stage renal disease, as the kidney graft provided a useful marker for rejection and often shared immunological responses with the pancreas. Further refinements in surgical techniques, such as the shift from bladder drainage to enteric drainage for exocrine secretions and improved organ preservation methods, continued to enhance the safety and efficacy of the procedure throughout the 1990s and into the 21st century, establishing pancreas transplant as a viable and life-changing treatment option.

Pancreas transplantation is a highly specialized procedure indicated for specific, severe manifestations of diabetes mellitus. The primary indications include:

\begin{itemize}
  \item Type 1 diabetes mellitus with end-stage renal disease (ESRD), typically performed as a simultaneous pancreas-kidney (SPK) transplant to address both conditions concurrently.
  \item Type 1 diabetes mellitus with recurrent, severe metabolic complications, such as frequent episodes of severe hypoglycemia requiring medical assistance or resulting in seizures or coma, despite optimal medical management.
  \item Type 1 diabetes mellitus with clinically significant hypoglycemia unawareness, where the patient loses the ability to perceive the warning signs of low blood sugar, leading to dangerous and unpredictable events.
  \item Extreme glycemic lability, characterized by wide and unpredictable fluctuations in blood glucose levels that are difficult to control with intensive insulin therapy and significantly impair quality of life.
  \item Progressive secondary complications of diabetes (e.g., severe neuropathy, advanced retinopathy, gastroparesis) that are significantly impacting quality of life and are expected to improve with stable euglycemia.
  \item Pancreas transplant after a prior kidney transplant (PAK) in patients with type 1 diabetes who have already received a kidney for diabetic nephropathy but continue to suffer from severe diabetic complications.
  \item Highly selected cases of type 2 diabetes mellitus, typically non-obese patients with minimal insulin resistance and severe complications, often in conjunction with a kidney transplant, though this remains controversial.
\end{itemize}

Pancreas transplantation can serve as both a primary and a secondary definitive treatment, depending on the clinical context and the patient's specific disease burden. For patients with type 1 diabetes and end-stage renal disease (ESRD), simultaneous pancreas-kidney (SPK) transplantation is often considered a primary, first-line definitive treatment. In this scenario, it addresses both the renal failure and the underlying diabetes with a single surgical event, offering the best long-term outcomes for both organs and significantly improving patient quality of life and survival compared to kidney transplant alone. The immunological benefits of transplanting both organs from the same donor also contribute to its primary role, as the kidney often serves as an immunological sentinel for rejection.

Conversely, pancreas transplant alone (PTA) or pancreas after kidney (PAK) transplant are typically considered secondary or salvage procedures. PTA is reserved for highly selected patients with type 1 diabetes who have severe, life-threatening diabetic complications (e.g., hypoglycemia unawareness, extreme glycemic lability) but preserved renal function. PAK is performed for patients who have already received a kidney transplant for diabetic nephropathy but continue to experience severe diabetic complications that are not adequately managed by insulin therapy. In these secondary roles, the procedure aims to improve quality of life and prevent further progression of diabetic complications, but the risk-benefit profile requires careful consideration due to the need for lifelong immunosuppression in patients who may not have immediate life-threatening organ failure beyond their diabetes.

\section*{Relevant Surgical Anatomy}
The pancreas is a retroperitoneal organ with both endocrine and exocrine functions, located posterior to the stomach, extending from the duodenum to the spleen. It is divided into the head, uncinate process, neck, body, and tail. For transplantation, the entire pancreas, along with a segment of the donor duodenum, is typically harvested to facilitate vascular anastomoses and exocrine drainage. The critical vascular supply for the donor pancreas includes the splenic artery (a branch of the celiac axis) and the superior mesenteric artery (a direct branch from the aorta), which are often anastomosed together via a Y-graft of donor iliac artery. Venous drainage is via the portal vein, formed by the confluence of the splenic and superior mesenteric veins.

In the recipient, the transplanted pancreas is typically placed in the right iliac fossa. The arterial Y-graft is anastomosed end-to-side to the recipient's common or external iliac artery, providing arterial inflow. The donor portal vein is anastomosed end-to-side to the recipient's common or external iliac vein for venous outflow. The donor duodenal segment, which contains the ampulla of Vater and allows for the drainage of pancreatic exocrine secretions, is anastomosed to the recipient's small intestine (typically the jejunum) for enteric drainage. The native pancreas is left in situ. The transplanted pancreas is denervated, meaning it no longer receives direct autonomic innervation, which can affect its physiological responses but does not impair its ability to produce insulin in response to glucose. Lymphatic drainage of the graft is also disrupted, which may contribute to a period of relative immune privilege immediately post-transplant.

\section{Surgical landmarks for recipient vessel preparation include:}
\begin{itemize}
  \item The common iliac artery and vein, or the external iliac artery and vein, in the chosen iliac fossa.
  \item The recipient's jejunum for enteric drainage, typically a Roux-en-Y limb or a simple loop.
\end{itemize}

\section*{Preoperative Workup \& Preparation}
Extensive pre-operative preparation is crucial for pancreas transplant recipients to ensure optimal outcomes and minimize risks. This comprehensive evaluation typically involves:

\begin{itemize}
  \item \textbf{Medical Evaluation:} A thorough assessment of cardiac function (ECG, echocardiogram, stress test, coronary angiography if indicated), pulmonary function (PFTs), and vascular status (Doppler studies, angiography) to ensure the patient can tolerate a major surgery and lifelong immunosuppression.
  \item \textbf{Infection Screening:} Comprehensive screening for active infections (viral, bacterial, fungal, parasitic), including HIV, hepatitis B and C, CMV, EBV, and tuberculosis, as active infections are absolute contraindications. Dental clearance is also required to eliminate potential sources of infection.
  \item \textbf{Immunological Workup:} Human Leukocyte Antigen (HLA) typing and crossmatching with potential donors to minimize the risk of rejection. Antibody screening is also performed.
  \item \textbf{Psychosocial Assessment:} Evaluation of the patient's understanding of the procedure, commitment to lifelong medication adherence, and availability of social support, which are critical for long-term success.
  \item \textbf{Nutritional Assessment:} Optimization of nutritional status, as malnutrition can impair healing and increase complication rates.
  \item \textbf{Medication Adjustments:} Insulin regimens are carefully managed up to the time of surgery. Other medications, such as anticoagulants, may need to be temporarily held or adjusted according to specific protocols.
  \item \textbf{NPO Status:} The patient must be nil per os (NPO) for a specified period before surgery to reduce the risk of aspiration during anesthesia.
  \item \textbf{Bowel Preparation:} If enteric drainage is planned, a mechanical bowel preparation may be administered to reduce the bacterial load in the colon and small intestine.
  \item \textbf{Pre-operative Antibiotics:} Prophylactic broad-spectrum antibiotics are administered intravenously just prior to incision to reduce the risk of surgical site infection.
  \item \textbf{Informed Consent:} Detailed discussion with the patient and family regarding the risks, benefits, alternatives, and the necessity of lifelong immunosuppression, ensuring full understanding and agreement.
  \item \textbf{Anesthesia Clearance:} A full pre-anesthetic evaluation to assess fitness for general anesthesia and identify any potential anesthetic risks.
\end{itemize}

Pancreas transplantation is a major surgical procedure with significant risks, necessitating careful patient selection to optimize outcomes.

\textbf{Absolute Contraindications:}
\begin{itemize}
  \item Active malignancy or a history of malignancy with a high risk of recurrence within a specified timeframe (e.g., 2-5 years, depending on cancer type).
  \item Active, uncontrolled infection (e.g., sepsis, osteomyelitis, active tuberculosis, severe dental infection) that could be exacerbated by immunosuppression.
  \item Severe, irreversible cardiac or pulmonary disease that would preclude safe surgery or long-term survival, such as severe coronary artery disease not amenable to revascularization or severe chronic obstructive pulmonary disease.
  \item Uncorrectable peripheral or visceral vascular disease that would compromise graft perfusion or make surgical anastomoses technically impossible.
  \item Morbid obesity (typically BMI > 30-35 kg/m\textasciicircum2, depending on center) due to increased surgical risk, technical challenges, and higher rates of graft complications.
  \item Active substance abuse or documented non-adherence to medical regimens, indicating an inability to comply with the rigorous demands of lifelong immunosuppression and follow-up.
  \item Severe, irreversible neurological damage or other conditions limiting rehabilitation potential and the ability to benefit from the transplant.
\end{itemize}

\textbf{Relative Contraindications and Precautions:}
\begin{itemize}
  \item Advanced age (typically > 60-65 years), as surgical risks and complications of immunosuppression increase, requiring careful individual assessment.
  \item Significant peripheral vascular disease or coronary artery disease that is not amenable to complete revascularization, increasing perioperative cardiovascular risk.
  \item Extensive prior abdominal surgery, which may complicate surgical access, increase operative time, and raise the risk of adhesions and bowel injury.
  \item Severe diabetic neuropathy, particularly autonomic neuropathy, which can increase perioperative risks such as gastroparesis, cardiovascular instability, and bladder dysfunction.
  \item Pancreas transplant alone (PTA) is considered with greater precaution due to the risks of lifelong immunosuppression in patients with otherwise preserved renal function, requiring a very high burden of diabetic complications to justify the procedure.
  \item Type 2 diabetes mellitus, due to higher rates of graft failure, recurrence of metabolic issues, and often significant comorbidities, requiring very strict selection criteria and careful patient counseling.
  \item Significant psychiatric illness that is not well-controlled, which may impact adherence to complex medical regimens.
\end{itemize}

\section*{Surgical Technique \& Procedure}
Pancreas transplantation is a complex surgical procedure performed under general anesthesia, typically lasting 4-6 hours, often longer if a kidney is transplanted simultaneously. The patient is positioned supine, and a lower midline abdominal incision or a transverse incision in the right or left iliac fossa is made.

The key operative steps include:

\begin{itemize}
  \item \textbf{Donor Organ Preparation (Bench Surgery):} While the recipient is being prepared, the donor pancreas, often with a segment of the duodenum, is meticulously prepared on a back table. This involves careful dissection of the vascular supply, including the splenic artery, superior mesenteric artery, and portal vein, and ensuring the integrity of the duodenal segment. A Y-graft of donor iliac artery is often used to connect the splenic and superior mesenteric arteries.
  \item \textbf{Recipient Vessel Preparation:} In the recipient, the common or external iliac artery and vein are identified and dissected in the chosen iliac fossa. These vessels will serve as the inflow and outflow for the transplanted pancreas, requiring careful mobilization and control.
  \item \textbf{Vascular Anastomoses:} The donor arterial Y-graft (connecting the splenic and superior mesenteric arteries) is anastomosed end-to-side to the recipient's common or external iliac artery. The donor portal vein is anastomosed end-to-side to the recipient's common or external iliac vein. This establishes robust arterial inflow and venous outflow for the new pancreas, crucial for graft viability.
  \item \textbf{Exocrine Drainage Management:} The donor duodenal segment is then anastomosed to the recipient's gastrointestinal tract, usually the jejunum (enteric drainage), or less commonly, to the bladder (bladder drainage). Enteric drainage is preferred to allow pancreatic enzymes to drain physiologically into the bowel, reducing metabolic complications and the risk of urinary tract issues.
  \item \textbf{Kidney Transplant (if SPK):} If a simultaneous pancreas-kidney transplant is performed, the donor kidney is typically implanted in the contralateral iliac fossa. Its renal artery is anastomosed to the recipient's iliac artery, and the renal vein to the iliac vein. The ureter is then anastomosed to the recipient's bladder.
  \item \textbf{Hemostasis and Drain Placement:} Meticulous hemostasis is achieved throughout the procedure to minimize blood loss. Surgical drains are often placed near the graft to monitor for fluid collections, pancreatic leaks, or bleeding.
  \item \textbf{Closure:} The abdominal incision is closed in layers, ensuring proper anatomical reconstruction and minimizing the risk of incisional hernia.
\end{itemize}
Throughout the procedure, careful attention is paid to maintaining organ viability, preventing ischemia-reperfusion injury, and minimizing blood loss. The surgical team works in close coordination with anesthesiologists to manage fluid balance, blood pressure, and glucose levels.

\section*{Complications \& Risks}
Pancreas transplantation carries significant risks, both surgical and related to lifelong immunosuppression. These can be broadly categorized as follows:

\textbf{General Surgical Risks:}
\begin{itemize}
  \item \textbf{Anesthesia Complications:} Reactions to medications, respiratory or cardiac arrest, stroke, pneumonia, and aspiration.
  \item \textbf{Bleeding:} Intraoperative or postoperative hemorrhage, potentially requiring blood transfusions or reoperation to control.
  \item \textbf{Infection:} Surgical site infection, pneumonia, urinary tract infection, or systemic sepsis, particularly concerning in immunosuppressed patients.
  \item \textbf{Thromboembolism:} Deep vein thrombosis (DVT) or pulmonary embolism (PE), which can be life-threatening.
  \item \textbf{Wound Complications:} Incisional hernia, wound dehiscence, or seroma formation.
\end{itemize}

\textbf{Pancreas Transplant-Specific Risks:}
\begin{itemize}
  \item \textbf{Graft Thrombosis:} Clot formation in the transplanted pancreas's arterial or venous vessels, leading to graft ischemia and failure, often requiring immediate reoperation and graft removal. This is a leading cause of early graft loss.
  \item \textbf{Pancreatic Leak/Fistula:} Leakage of pancreatic enzymes or intestinal contents from the anastomotic sites (duodenojejunostomy), which can lead to severe peritonitis, abscess formation, or skin fistulas, requiring prolonged drainage or reoperation.
  \item \textbf{Graft Pancreatitis:} Inflammation of the donor pancreas, often due to ischemia-reperfusion injury or rejection, which can impair graft function and cause abdominal pain.
  \item \textbf{Rejection:} The recipient's immune system attacking the transplanted organ, which can be acute (occurring days to months post-transplant) or chronic (gradual loss of function over years). Requires increased immunosuppression and can lead to graft failure.
  \item \textbf{Infection related to Immunosuppression:} Increased susceptibility to opportunistic infections (e.g., Cytomegalovirus (CMV), fungal infections, Pneumocystis jirovecii pneumonia (PJP), Epstein-Barr virus (EBV)) due to the weakened immune system.
  \item \textbf{Complications of Immunosuppression:}
    \begin{itemize}
      \item Nephrotoxicity (kidney damage), particularly with calcineurin inhibitors like tacrolimus and cyclosporine.
      \item Hypertension, dyslipidemia, and new-onset diabetes after transplant (NODAT, if not already diabetic or if the transplanted pancreas fails).
      \item Increased risk of certain malignancies, especially skin cancers and post-transplant lymphoprolroliferative disorder (PTLD).
      \item Bone marrow suppression, gastrointestinal upset, tremors, neurotoxicity, and osteoporosis.
    \end{itemize}
  \item \textbf{Anastomotic Stricture or Obstruction:} Narrowing or blockage at the site where the donor duodenum is connected to the recipient's bowel, potentially causing bowel obstruction or impaired exocrine drainage.
  \item \textbf{Cyst Formation:} Pseudocyst or true cyst formation around the graft.
  \item \textbf{Gastrointestinal Complications:} Diarrhea, malabsorption, or peptic ulcer disease, sometimes related to immunosuppression or altered anatomy.
\end{itemize}

\section*{Postoperative Care \& Recovery}
Immediately following a pancreas transplant, the patient is transferred to an intensive care unit (ICU) for close monitoring. This critical period involves continuous assessment of vital signs, blood glucose levels, urine output, and drain output. Intravenous fluids, pain medication, and immunosuppressive drugs are administered. Blood glucose levels are meticulously managed, often with an insulin drip initially, until the transplanted pancreas begins to function consistently, typically within hours to days. Frequent blood tests monitor pancreatic enzymes (amylase, lipase), kidney function (creatinine), and immunosuppressant drug levels.

Over the first 24-48 hours, if stable, the patient may be transferred to a specialized transplant ward. Oral intake is gradually advanced from clear liquids to a regular diet as bowel function returns. Early ambulation is strongly encouraged to prevent complications like deep vein thrombosis and pneumonia. The surgical drains are typically removed when output is minimal and non-suspicious. The patient and family receive extensive education on their new medication regimen, including the critical importance of lifelong immunosuppression, potential side effects, and warning signs of complications. Discharge usually occurs within 1-3 weeks, depending on recovery progress, the absence of significant complications, and the patient's ability to manage their medications and care. Long-term, patients require lifelong adherence to immunosuppressive medications and regular follow-up appointments to monitor graft function, detect rejection early, and manage potential side effects of immunosuppression.

During a pancreas transplant, the surgeon documents several key findings. These include the appearance of the donor pancreas (e.g., healthy, well-perfused, no obvious injury), the quality of the vascular anastomoses (e.g., patent, good flow, no kinks), and the integrity of the duodenojejunostomy. Any fluid collections, signs of inflammation, or unexpected anatomical variations are also noted. The successful establishment of blood flow to the graft is often confirmed visually by the pinking up of the pancreas and duodenum, and sometimes by Doppler ultrasound intraoperatively.

The pathology report on resected tissues, typically obtained from a graft biopsy performed if rejection is suspected, provides crucial diagnostic information. The pathologist examines tissue samples for signs of acute cellular rejection (e.g., lymphocytic infiltration of the exocrine parenchyma and/or islets, arteritis), antibody-mediated rejection (e.g., C4d staining in peritubular capillaries, neutrophilic capillaritis), or other pathologies such as ischemia-reperfusion injury, pancreatitis, or infection. The report details the severity and type of rejection, guiding subsequent adjustments to immunosuppressive therapy. In cases of graft pancreatectomy due to irreversible failure, the entire explanted pancreas is sent for detailed pathological examination to determine the cause of failure.

A normal, successful postoperative course following pancreas transplantation is characterized by a progressive return to health and the achievement of insulin independence. Patients typically experience a gradual resolution of surgical pain, which is managed with analgesics and progressively weaned. Within days to a week, the transplanted pancreas usually begins to function, leading to a normalization of blood glucose levels and the cessation of exogenous insulin requirements. Patients will notice a significant improvement in their energy levels and overall well-being as their metabolic control is restored.

Wound healing progresses normally, with sutures or staples typically removed around 2-3 weeks post-surgery. Patients are encouraged to gradually increase their activity levels, returning to light activities within a few weeks and more strenuous activities over several months. The fear of hypoglycemic episodes, a major burden for many diabetic patients, is eliminated. The long-term expectation is stable graft function, freedom from insulin injections, and the prevention or reversal of diabetic complications, allowing for a significantly improved quality of life, provided there is strict adherence to the lifelong immunosuppressive regimen and regular medical follow-up.

Lifelong follow-up is essential after a pancreas transplant to ensure graft function, manage immunosuppression, and monitor for complications.

\begin{itemize}
  \item \textbf{Regular Clinic Visits:} Frequent visits initially (weekly, then monthly) gradually decreasing to quarterly or semi-annually, depending on stability and center protocols.
  \item \textbf{Blood Tests:}
    \begin{itemize}
      \item Blood glucose, HbA1c, and C-peptide to monitor graft endocrine function and confirm insulin independence.
      \item Serum amylase and lipase to monitor graft exocrine function and detect pancreatitis or leaks.
      \item Immunosuppressant drug levels (e.g., tacrolimus, cyclosporine, sirolimus) to ensure therapeutic ranges and minimize toxicity.
      \item Complete blood count (CBC), liver function tests (LFTs), kidney function tests (creatinine, GFR) to monitor for side effects of immunosuppression and overall health.
      \item Viral load monitoring (e.g., CMV, EBV, BK virus) due to increased infection risk from immunosuppression.
    \end{itemize}
  \item \textbf{Imaging Studies:} Regular Doppler ultrasound of the transplanted pancreas to assess blood flow, rule out vascular complications (thrombosis, stenosis), and detect fluid collections.
  \item \textbf{Graft Biopsy:} Performed if there is suspicion of rejection (e.g., rising glucose, enzymes, or creatinine) to confirm diagnosis and guide treatment adjustments.
  \item \textbf{Dietary Counseling:} Guidance on healthy eating to maintain a healthy weight, manage potential metabolic side effects of immunosuppression (e.g., dyslipidemia, hypertension), and prevent weight gain.
  \item \textbf{Physical Therapy/Rehabilitation:} As needed, to regain strength and mobility post-surgery and improve overall physical conditioning.
  \item \textbf{Screening for Immunosuppression Complications:} Regular screening for skin cancer, cardiovascular disease, osteoporosis, and other malignancies, which are increased risks with long-term immunosuppression.
  \item \textbf{Vaccinations:} Adherence to recommended vaccination schedules, with careful consideration for live vaccines, which are generally contraindicated in immunosuppressed patients.
\end{itemize}
Patients are educated to immediately report warning signs such as fever, new or worsening abdominal pain, elevated blood glucose, changes in urine output, or general malaise to their transplant team, as these could indicate serious complications like infection or rejection.

\section*{Outcomes \& Prognosis}
Pancreas transplantation is a relatively rare procedure compared to kidney or liver transplantation, with approximately 1,000 to 1,500 procedures performed annually worldwide, predominantly in North America and Europe. The majority of recipients are individuals with type 1 diabetes mellitus, often complicated by end-stage renal disease (ESRD). Simultaneous pancreas-kidney (SPK) transplants account for the largest proportion, typically around 70-80\% of all pancreas transplants, reflecting the strong indication for patients with diabetic nephropathy. Pancreas transplant alone (PTA) and pancreas after kidney (PAK) transplants are less common, reserved for highly selected patients with severe diabetic complications but without concurrent ESRD or after a prior kidney transplant, respectively. Recipients are typically younger adults, often in their 30s to 50s, with no significant sex predilection. While the incidence of type 1 diabetes is relatively stable, the number of eligible candidates for pancreas transplant is limited by strict selection criteria, the complexity of the procedure, and donor organ availability. Mortality rates have significantly decreased over the decades due to improved surgical techniques and immunosuppression, with 1-year patient survival rates now exceeding 95\% for SPK recipients, though morbidity remains substantial due to the complexity of the surgery and the lifelong need for immunosuppression.

The outcomes and prognosis following pancreas transplantation have dramatically improved, offering significant benefits for carefully selected patients. For simultaneous pancreas-kidney (SPK) recipients, 1-year patient survival rates typically exceed 95\%, with 5-year survival rates around 85-90\%. Pancreas graft survival rates (defined as insulin independence) are also excellent, with 1-year rates often above 85\% and 5-year rates ranging from 60-75\%. For pancreas transplant alone (PTA) and pancreas after kidney (PAK) recipients, patient survival is similar, but pancreas graft survival rates are slightly lower, reflecting the increased immunological challenge without a co-transplanted kidney and often a higher baseline risk profile.

The primary functional outcome is sustained insulin independence, which is achieved in the vast majority of successful transplants. This leads to a profound improvement in quality of life, freedom from the daily burden of insulin injections and glucose monitoring, and elimination of severe hypoglycemic episodes. Furthermore, successful pancreas transplantation can halt or even reverse the progression of many diabetic complications, including neuropathy, retinopathy, and gastroparesis, and significantly improve cardiovascular risk factors. Prognostic factors influencing outcomes include donor age, cold ischemia time, HLA matching, and recipient adherence to immunosuppressive therapy. While diabetes itself does not recur in the transplanted pancreas (as it is an autoimmune disease affecting the original pancreas), chronic rejection or complications of immunosuppression can lead to graft failure over time. Landmark studies and registries, such as those from the International Pancreas Transplant Registry (IPTR) and the Organ Procurement and Transplantation Network (OPTN)/United Network for Organ Sharing (UNOS), consistently demonstrate these favorable outcomes.

\section*{References}
\begin{enumerate}
  \item MedlinePlus. Pancreas transplant. U.S. National Library of Medicine; 2024. Available from: \url{https://medlineplus.gov/pancreastransplant.html}
  \item Sabiston Textbook of Surgery: The Biological Basis of Modern Surgical Practice. 22nd ed. Elsevier; 2022.
  \item Schwartz's Principles of Surgery. 11th ed. McGraw-Hill Education; 2019.
  \item American Society of Transplant Surgeons (ASTS). Pancreas Transplantation. Available from: \url{https://asts.org/patients/organ-transplantation/pancreas-transplantation}
\end{enumerate}

\clearpage
